\chapter{Análise e Discussão}\label{cap:discussao}

Este capítulo interpreta os resultados do \autoref{cap:resultados} e os situa em relação à
literatura. A discussão organiza-se em torno de três pontos: a leitura dos campos de coerência e de
potência obtidos, a comparação do método executado com os trabalhos anteriores no mesmo arranjo e na
análise espectral de padrões de escoamento, e as limitações que delimitam o alcance das conclusões.

% =========================================================
\section{Interpretação dos campos obtidos}\label{sec:disc-interpretacao}
% =========================================================

A coerência \textit{Wavelet} alternativa foi adotada como descritor principal, na expectativa de que
o acoplamento espectral entre dois ensaios de um mesmo padrão se manifestasse como uma banda de
coerência elevada nos comprimentos de onda característicos daquele padrão. Os campos obtidos, porém,
apresentaram coerência baixa e irregular no atraso nulo, sem uma banda estável que acompanhasse a
faixa dominante do espectro de potência. Uma leitura consistente com esse comportamento é a de que
dois ensaios distintos de um mesmo padrão, ainda que reproduzam o mesmo tipo de estrutura de
escoamento, não compartilham uma fase espacial comum: por serem realizações turbulentas
independentes, coletadas em instantes e condições microscopicamente distintas, o alinhamento entre
as estruturas de um ensaio e as do outro não se mantém ao longo do percurso. Como a coerência mede a
constância da relação de fase entre os dois sinais em cada escala, a ausência de uma fase espacial
comum reduz a coerência avaliada no atraso nulo, mesmo quando os dois ensaios têm energia espectral
concentrada na mesma faixa de comprimento de onda.

O perfil de potência média $|\text{WPS}|$ descreveu os padrões de forma mais consistente justamente
porque não depende da relação de fase entre os ensaios. O espectro de potência resume a distribuição
de energia de cada ensaio por comprimento de onda, uma grandeza que se mostrou reprodutível dentro de
cada padrão. Essa observação alinha-se ao uso de descritores espectrais escalares por ensaio na
literatura de identificação de padrões de escoamento. \citeonline{nguyen2010} aplicaram a
transformada \textit{Wavelet} contínua a sinais de fração de vazio e extraíram o coeficiente de
energia \textit{Wavelet} local e a escala efetiva como critérios objetivos de discriminação dos
padrões, e \citeonline{sun2012} identificaram padrões gás-líquido pela entropia dos coeficientes da
decomposição \textit{Wavelet} das flutuações de pressão. Em ambos os casos, o descritor que separa os
padrões deriva da distribuição de energia espectral do próprio sinal, e não de uma medida de fase
relativa entre sinais, o que é coerente com o desempenho comparativamente melhor do perfil de
potência média observado neste trabalho. A hipótese de estacionariedade sobre a qual a formulação de
\citeonline{schuck2019} se apoia permanece necessária tanto para o espectro de potência quanto para a
coerência; a diferença de desempenho entre os dois descritores diz respeito à dependência de fase, e
não à validade da hipótese.

% =========================================================
\section{Comparação com a literatura}\label{sec:disc-literatura}
% =========================================================

O trabalho de \citeonline{deoliveira2021}, conduzido no mesmo arranjo experimental do LETEF,
caracterizou os padrões de escoamento aplicando a transformada \textit{Wavelet} contínua aos sinais
dos sensores no domínio do tempo e extraindo as frequências características de cada padrão. A
abordagem executada neste trabalho difere daquela em dois aspectos: o domínio de análise, que passou
do tempo para o espaço, com o sinal reconstruído do panorama das imagens em vez dos sinais dos
sensores; e o descritor, que passou do espectro de potência de um sinal individual para a coerência
alternativa entre pares de ensaios. A parte do método que mais se aproximou dos resultados daquele
trabalho foi justamente o perfil de potência média, que, assim como a caracterização de
\citeonline{deoliveira2021}, resume a energia espectral de cada ensaio, agora em função do
comprimento de onda espacial.

Em relação à formulação clássica da coerência \textit{Wavelet} \cite{torrence1998,grinsted2004}, a
formulação alternativa de \citeonline{schuck2019} dispensa a suavização no plano tempo-escala. Na
definição clássica, a coerência é construída a partir de espectros cruzados suavizados, e a
suavização é o que impede que a medida seja trivialmente unitária, funcionando como um coeficiente de
correlação localizado no plano tempo-frequência \cite{grinsted2004}. A formulação alternativa
substitui essa suavização pelo cálculo sobre as funções de correlação avaliadas no atraso nulo. A
aplicação dessa formulação a sinais espaciais de brilho, no contexto de identificação de padrões de
escoamento bifásico, não havia sido reportada, e os resultados aqui obtidos indicam que o
comportamento da coerência alternativa entre ensaios distintos merece cautela: o descritor foi
sensível à ausência de fase espacial comum entre realizações independentes, ponto que a literatura
clássica trata pela suavização e que a variante de \citeonline{schuck2019} trata de outra forma.

% =========================================================
\section{Limitações}\label{sec:disc-limitacoes}
% =========================================================

Três limitações delimitam o alcance dos resultados. A primeira é o número reduzido de ensaios por
padrão: apenas \textit{churn} e anular ondulado dispõem de ensaios suficientes para sustentar uma
leitura de repetibilidade, enquanto \textit{plug}, \textit{slug} e \textit{slug}--\textit{churn} têm
poucos ensaios e foram examinados de forma ilustrativa. A segunda é a possível influência das
condições operacionais: como cada ensaio foi coletado sob uma combinação específica de fluxo mássico
e fluxo de calor, a variação da assinatura espectral entre padrões pode estar parcialmente associada
à variação dessas condições, e a separação entre o efeito do padrão e o efeito das condições
operacionais não foi isolada. A terceira é o desempenho da coerência alternativa aquém do esperado
como descritor discriminante, discutido na \autoref{sec:disc-interpretacao}, que deslocou o peso da
caracterização para o perfil de potência média. A quantificação estatística formal dessas leituras,
por fim, não foi executada e permanece como encaminhamento no \autoref{cap:conclusao}.
